\chapter{How to review the code stage by stage}

The repository is intentionally organized around the six project stages. Focus on one stage and its tests before moving downstream.

\section{Stage 1: data}

\textbf{Question to answer:} Can every training/evaluation record be traced to valid geometry, the correct split, and a verified source release?

Start with:

\begin{itemize}[leftmargin=6mm]
  \item \file{main\_project/src/consentguard/stage\_01\_data/dataset.py};
  \item \file{main\_project/src/consentguard/stage\_01\_data/data\_quality.py};
  \item preprocessing/validation scripts under the Stage 1 scripts directory; and
  \item Stage 1 tests for empty targets, geometry, taxonomy, and same-release manifests.
\end{itemize}

Review order: record schema, image loading, positive target construction, empty target construction, resize transforms, geometry checks, manifest provenance, and split-leakage checks.

\section{Stage 2: baseline model}

\textbf{Question to answer:} Can a configuration deterministically construct, train, resume, and evaluate the intended model?

The training code is segmented into:

\begin{itemize}[leftmargin=6mm]
  \item \file{config.py}: validated experiment settings;
  \item \file{models.py}: model registry and custom heads;
  \item \file{data\_loading.py}: dataset/sampler/dataloader construction;
  \item \file{optimization.py}: optimizer and scheduler;
  \item \file{training\_loop.py}: epochs, evaluation, checkpointing, resume;
  \item \file{metrics.py}: COCO-style evaluation; and
  \item \file{diagnostics.py}: coverage, leakage, over-redaction, and proposals.
\end{itemize}

Read tests before changing behavior. They define expected config validation, checkpoint compatibility, class-agnostic mask behavior, atomic artifact writes, and privacy metrics.

\section{Stage 3: specialists}

\textbf{Question to answer:} Does each external tool produce correct typed evidence or an explicit unavailable/error state?

Review \file{common.py}, then each provider adapter, then \file{orchestrator.py}. Confirm:

\begin{itemize}[leftmargin=6mm]
  \item provider names are unique and stable;
  \item scores are bounded and geometry clipped;
  \item versions and source detection IDs are retained;
  \item recognized OCR strings are not stored unnecessarily;
  \item missing weights/dependencies raise provider-unavailable errors; and
  \item execution order is deterministic and GPU-safe.
\end{itemize}

\section{Stage 4: fusion and calibration}

\textbf{Question to answer:} Can evidence from different frameworks be compared, transformed, fused, and explained without losing origin?

Review domain models, provider base contracts, geometry helpers, threshold registry, and fusion logic. Use Stage 4 tests to understand exact/fallback rules, mask RLE round-trips, overlap grouping, provenance retention, rejected evidence, and provider state.

The candidate YAML is not release-ready. Do not confuse readable configuration with validated calibration.

\section{Stage 5: review and export}

\textbf{Question to answer:} Does the approved human decision become exactly the fresh output pixels, and does any uncertainty block export?

Trace the pipeline through normalizer, manual-mask editor, prediction/oracle renderers, policy engine, and assurance service. Inspect tests covering source-overwrite refusal, fresh decode, per-class thresholds, fail-closed consent/policy, independent attackers, and manual layer union.

\section{Stage 6: release}

\textbf{Question to answer:} Can one machine-readable report prove every gate passed?

The release-gate module already encodes baseline mAP, overall recall/leakage, negative false positives, three seeds, fail-closed assurance, domain recall, and per-class recall. The missing work is producing honest metrics from frozen target evaluations, not adding a button that says release.

\section{Suggested review checklist}

For every file or change, ask:

\begin{enumerate}[leftmargin=7mm]
  \item What input schema does this code trust?
  \item What output or state does it promise?
  \item How does it fail, and is failure distinguishable from an empty result?
  \item Which coordinate system or class ontology is active?
  \item Which configuration and version affect behavior?
  \item What test proves the important contract?
  \item Could a stale artifact or hidden default change the result?
  \item Does the code reveal or unnecessarily retain sensitive information?
\end{enumerate}

\section{Commands to understand, not memorize}

Use the project environment and the Stage-specific scripts. Before any long run, execute a data validation, environment preflight, loader smoke test, and one-batch learning smoke test. Verify the resolved config and output directory before starting Kaggle.

The exact commands may evolve; the invariant is the sequence:

\begin{lstlisting}[language={}]
validate data -> preflight environment -> loader smoke test
-> learning smoke test -> full train/resume -> evaluate
-> diagnostics -> compare against frozen baseline
\end{lstlisting}

Never run the locked test command during development.

